\chapter{System Design and Architecture}

\section{System Overview}

\textit{DEX} is designed as a distributed framework intended to provide a
unified logical view over heterogeneous and distributed data sources, accessed
strictly in read-only mode. The framework adopts a federated query processing
strategy: queries are formulated over the logical model defined above the
concrete schemas of the underlying databases, providing a layer of abstraction
that hides how data is physically stored and structured. Those queries are then
decomposed into subqueries executed against the underlying sources, and
subsequently resolved by aggregating partial results. By structuring
\textit{DEX} as a distributed system, subquery execution is carried out in
parallel, thereby maximizing throughput and maintaining system responsiveness. 

User interaction with \textit{DEX} is mediated by a graphical interface that
enables the construction of queries over the logical model through visual
filters, eliminating the need for explicit query coding. Upon formulation, the
system translates the global query into localized subqueries and orchestrates
their execution. Each subquery is assigned to an individual worker process
targeting a specific data source, while partial results converge at a central
logical point for final aggregation once all responses are received.
Furthermore, secondary reduction operations, such as summations, averages, and
counts, can be applied directly to the aggregated dataset. 

To mitigate the computational overhead associated with unexpected failures, the
design incorporates systemic fault tolerance alongside caching mechanisms.
Operating under the premise that network disruptions and component crashes are
inherent to distributed environments, the architecture is structured to preserve
partial execution progress and isolate recovery strictly to failing entities,
avoiding when is possible the redundant re-execution of the global workload,
optimizing resource utilization and sustaining operational efficiency.  

Ultimately, this design encapsulates the complexities of data distribution,
source heterogeneity, execution orchestration and fault recovery. From the user
perspective, the entire execution pipeline is exposed as a single, seamless
operation abstracting the underlying distributed execution offered by
\textit{DEX}. 
